\section{Counting Towers}
\label{sec:cses-counting-towers}

\problemstatement{Build a tower of width $2$ and height $n$ out of blocks
that are $1\times1$ or $1\times2$ (blocks may be placed in either
orientation to fill the $2$-wide grid). Count the number of distinct
towers of height $n$, modulo $10^9+7$. Answer $T$ test cases.}

\begin{patternbox}
This is a row-by-row DP with a two-valued state describing the
\textbf{boundary shape} between consecutive levels. At each height the
seam between the level below and above is either \emph{split} (two
separate cells) or \emph{joined} (one $1\times2$ block spanning the
width). Transitions between these two boundary states, counted with fixed
multipliers, give the recurrence.
\end{patternbox}

\subsection*{Two boundary states}

Process the tower one unit of height at a time, tracking the shape of the
horizontal seam at the current top:
\begin{itemize}
  \item State $0$ -- the seam is \textbf{split} into two independent
        cells (the two columns are separable at this height).
  \item State $1$ -- the seam is \textbf{joined} by a single horizontal
        $1\times2$ block spanning both columns.
\end{itemize}
Counting how each state can be extended upward by one level yields:
\[
  \begin{aligned}
    \textit{dp}[i][0] &= 4\,\textit{dp}[i-1][0] + 2\,\textit{dp}[i-1][1],\\
    \textit{dp}[i][1] &= \textit{dp}[i-1][0] + \textit{dp}[i-1][1].
  \end{aligned}
\]
The constants $4$ and $2$ enumerate the ways to fill one more level given
the seam below (four fillings keep the seam split, etc.). Base:
$\textit{dp}[1][0]=\textit{dp}[1][1]=1$.

\begin{ideabox}[The State Is the Seam, Not the Whole Tower]
The only information the next level needs from everything below is the
\emph{shape of the current top seam}. Compressing the entire history into
one of two boundary states is what makes this an $O(n)$ DP rather than an
exponential enumeration -- a miniature ``broken profile'' idea, revisited
fully in Counting Tilings (Section~\ref{sec:cses-counting-tilings}).
\end{ideabox}

\subsection*{Algorithm}

\begin{enumerate}
  \item Precompute $\textit{dp}[i][0],\textit{dp}[i][1]$ for
        $i=1\ldots\max n$ using the recurrence above.
  \item For each test case $n$, answer
        $\textit{dp}[n][0]+\textit{dp}[n][1]$ modulo $10^9+7$.
\end{enumerate}

\subsection*{Dry run}

\dryrun{Small heights.}

\begin{itemize}
  \item $n=1$: $\textit{dp}[1][0]=1,\textit{dp}[1][1]=1$; total $2$
        (two cells, or one horizontal block).
  \item $n=2$: $\textit{dp}[2][0]=4\cdot1+2\cdot1=6$,
        $\textit{dp}[2][1]=1+1=2$; total $8$.
  \item $n=3$: $\textit{dp}[3][0]=4\cdot6+2\cdot2=28$,
        $\textit{dp}[3][1]=6+2=8$; total $36$.
\end{itemize}

\subsection*{C++ implementation}

\begin{lstlisting}
#include <bits/stdc++.h>
using namespace std;
const int MOD = 1e9 + 7;
const int MAXN = 1e6 + 1;

int main() {
    static long long dp0[MAXN], dp1[MAXN];
    dp0[1] = dp1[1] = 1;
    for (int i = 2; i < MAXN; i++) {
        dp0[i] = (4 * dp0[i-1] + 2 * dp1[i-1]) % MOD;
        dp1[i] = (dp0[i-1] + dp1[i-1]) % MOD;
    }

    int t;
    cin >> t;
    while (t--) {
        int n;
        cin >> n;
        cout << (dp0[n] + dp1[n]) % MOD << "\n";
    }
    return 0;
}
\end{lstlisting}

\begin{complexitybox}
\textbf{Time Complexity.} $O(\max n)$ precomputation, then $O(1)$ per
query. \textbf{Space Complexity.} $O(\max n)$ for the two DP arrays.
\end{complexitybox}

\begin{takeawaybox}
When a structure is built layer by layer, the DP state need only capture
the \emph{interface} between the built and unbuilt parts. Here two seam
states suffice; identifying that minimal boundary description is the core
skill, scaling up to bitmask profiles for wider grids.
\end{takeawaybox}
